\chapter{OPD 的目标、残差与访问测度}
\label{chap:opd}
\section{固定 prefix 的镜像问题}
令 $\nu$ 为固定 prefix 测度，$b$ 为行为策略，$T$ 为教师。在有限词表上假设
$b(a\mid s),T(a\mid s)>0$，定义停止梯度的信号
\[
 r_b(s,a)=\log T(a\mid s)-\log b(a\mid s),\qquad
 \Psi_{b,\eta}(p;s)=\langle p,r_b\rangle-\eta^{-1}\KL(p\Vert b).
\]
镜像温度为 $\tau=1/\eta$；参数 optimizer learning rate 另记为 $\gamma$。
\begin{theorem}[变分恒等式与唯一目标]\label{thm:gibbs-opd}
设 $\eta>0$，并定义
\[
 G_\eta(b)(a\mid s)=\frac{b(a\mid s)^{1-\eta}T(a\mid s)^\eta}{Z_b(s,\eta)},
 \quad Z_b(s,\eta)=\sum_a b(a\mid s)^{1-\eta}T(a\mid s)^\eta.
\]
则
\[
 \Psi_{b,\eta}(p;s)=\eta^{-1}\log Z_b(s,\eta)
                 -\eta^{-1}\KL(p\Vert G_\eta(b)).
\]
唯一最优解为 $G_\eta(b)$。
\end{theorem}
\begin{proof}
逐项取对数得 $\log G_\eta(b)=\log b+\eta r_b-\log Z_b$。
因此
\begin{align*}
 \KL(p\Vert G_\eta(b))
 &=\sum_a p(a)[\log p(a)-\log b(a)-\eta r_b(a)+\log Z_b]\\
 &=\KL(p\Vert b)-\eta\langle p,r_b\rangle+\log Z_b.
\end{align*}
移项得到恒等式。由 $-\log u\ge1-u$，取 $u=G_\eta(b)(a)/p(a)$
并对 $p$ 求和，得到 KL 非负性。等号要求 $p=G_\eta(b)$。
\end{proof}
$\eta=1$ 时 $\Psi_{b,1}(p)=-\KL(p\Vert T)$，anchor 项精确抵消。
$\eta<1$ 保留 anchor；$\eta>1$ 在 log-prob 坐标外推。
$b(a)=0$ 时 KL 近端禁止 $p(a)>0$；$T(a)=0$ 时信号负无穷。
这两类 support 问题需要单独处理，不能由无限 inner budget 消除。

\section{中心化残差的零点}
定义
\[
 h_b=r_b-\E_b r_b=r_b+\KL(b\Vert T),\qquad
 \delta_p=\eta h_b-\log(p/b)-\KL(b\Vert p).
\]
因为 $\E_b\log(p/b)=-\KL(b\Vert p)$，所以 $\E_b\delta_p=0$。
回归目标为
\[
 L_{\mathrm{res}}(p)=\E_{s\sim\nu,a\sim b}[\delta_p(s,a)^2].
\]
\begin{theorem}[残差与镜像目标的等价零点]\label{thm:residual}
若所有受训 prefix 上 $b$ 为正且策略族可表示 $G_\eta(b)$，
$L_{\mathrm{res}}=0$ 当且仅当 $p=G_\eta(b)$，以 $\nu$-几乎处处为准。
\end{theorem}
\begin{proof}
令 $q=G_\eta(b)$，则
\[
 \log(q/b)=\eta r_b-\log Z_b,\qquad
 \KL(b\Vert q)=-\eta\E_b r_b+\log Z_b.
\]
代入得
\[
 \delta_q=\eta(r_b-\E_b r_b)-(\eta r_b-\log Z_b)
          -(-\eta\E_b r_b+\log Z_b)=0.
\]
反向，非负平方的积分为零意味着 $\delta_p=0$。于是
$\log(p/b)=\eta h_b-c(s)$，其中 $c(s)=\KL(b\Vert p)$。
指数化并归一化，得到 $p\propto b e^{\eta h_b}=b e^{\eta r_b}$。
\end{proof}
同一零点不意味着同一投影。有限策略族中，reverse KL、forward KL 的加权最大似然、
平方 log-ratio 残差会产生不同最优参数；非凸驻点也不能由零点定理排除。
PMD-mean 用行为均值替代配分函数的具体损失，会引入隐式正则化和
Lambert-$W$ 驻点结构，不能概括为“所有中心化目标都与 exact PMD 同解”。
BPO 的 Bellman 残差还要保留未来价值与终局信号。

\section{序列级与逐状态级镜像}
序列级镜像把整个响应 $y$ 当成动作：
\[
 P^+(y\mid x)=B(y\mid x)e^{\eta R(x,y)}/Z(x).
\]
在确定性自回归树上，定义 continuation partition function
\[
 Z(s)=\E_{y_{\ge t}\sim B(\cdot\mid s)}[e^{\eta R(x,y)}],
 \quad Z(s_{\mathrm{terminal}})=e^{\eta R(x,y)}.
\]
条件期望给出 $Z(s)=\sum_a b(a\mid s)Z(sa)$。故
\[
 p^+(a\mid s)=b(a\mid s)\frac{Z(sa)}{Z(s)}.
\]
其概率和为一，且沿序列相乘得到
\[
 \prod_t\frac{p^+(a_t\mid s_t)}{b(a_t\mid s_t)}
 =\prod_t\frac{Z(s_{t+1})}{Z(s_t)}
 =\frac{e^{\eta R(x,y)}}{Z(x)}.
\]
这里包含后缀软价值。直接在每个 prefix 用即时 $\log T-\log b$ 做几何插值
一般不同于序列几何插值；即使 $R=\log(T(y)/B(y))$，后者的条件策略仍含
continuation normalization。
逐状态 PMD 用真实 $A^b(s,a)$ 时具有政策改进语义；
teacher log-ratio 没有自动成为环境优势。

\section{历史、动作与未来比率}
令 $d^\pi_t$ 为第 $t$ 步的 prefix 分布，
$d^\pi=\sum_t d^\pi_t$ 为未归一化 occupancy，$\ell_\pi=\E_\pi|y|$ 为其总质量。
对当前 integrand $f(s_t,a_t)$，目标密度为 $d^\pi_t(s)\pi(a\mid s)$。
行为数据为 $d^b_t(s)b(a\mid s)$ 时需要
\[
 \Omega_t(s)\rho_t(s,a)=\frac{d^\pi_t(s)}{d^b_t(s)}
                       \frac{\pi(a\mid s)}{b(a\mid s)}.
\]
自回归树历史唯一时 $\Omega_t=\prod_{i<t}\rho_i$。
完整词表 expectation 已经消除 action Monte Carlo，不再需要 action ratio。
平方残差的目标本来在固定 $(\nu,b)$ 下，故不需要 student sequence IS，
但它不自动成为当前 on-policy sequence objective 的无偏估计。

\begin{theorem}[历史校正的条件方差优势]
假设所有比率存在且平方可积，$W=\prod_{i=1}^{H}\rho_i$，
$W_{\le t}=\prod_{i\le t}\rho_i$，$f_t$ 对 $(s_t,a_t)$ 可测，则
\[
 \E_b[Wf_t\mid s_t,a_t]=W_{\le t}f_t,\qquad
 \Var_b(W_{\le t}f_t)\le\Var_b(Wf_t).
\]
\end{theorem}
\begin{proof}
目标与行为使用相同环境转移。给定历史与当前动作，未来比率对行为后缀积分为一，
因为目标后缀概率归一化；$W_{\le t}$ 在自回归树上对当前历史可测。
所以第一式成立。总方差公式
\[
 \Var(X)=\Var(\E[X\mid\mathcal G])+\E\Var(X\mid\mathcal G)
\]
取 $X=Wf_t$ 和 $\mathcal G=\sigma(s_t,a_t)$ 即得第二式。
\end{proof}
因此完整无偏 suffix ratio 并不会必然改变纯 prefix integrand 的期望，
但会加入不必要的条件方差。重复校正、截断或改变停止梯度语义才会引入额外偏差。
短窗口 IS 是偏差与方差折中，不能用“无序列 IS”来宣称零偏差。
若 $|\log\rho_i|\le B$，则 $W$ 最坏可达 $e^{HB}$，而
$\log W=\sum_i\log\rho_i$ 的幅度至多 $HB$；
这解释了残差回归的数值优势，但仍须控制 score 与样本二阶矩。

\section{覆盖反例与一步 teacher 投影}
考虑两步树，第一步为 $L,R$，训练数据只包含 $L$ 分支。
受训 prefix 的 teacher KL 可以降为零，而 $R$ 分支条件策略任意。
当第一步更新令 $R$ 更常被访问，第二轮 rollout 才暴露旧 batch 未约束的误差。
所以 $\eta=1$ 只在完整覆盖、可实现性、正确 logits 和全局 inner 求解同时成立时，
才能推出全局 $\pi=T$。
有限参数情况下正确极限是
\[
 \theta^\star\in\argmin_\theta\E_{s\sim\nu}
          \KL(\pi_\theta(\cdot\mid s)\Vert T(\cdot\mid s)).
\]
重新 rollout 更新的是访问测度，并不能用无限次旧 batch 优化代替。

